\section{Quantitative homological perturbation}
\label{sec:homological}

The effective nerve arguments in the existing packets obtain a small global
model when all relevant intersections are acyclic or carry supplied
contractions.  A useful extension retains a small model of each intersection
even when its homology is nontrivial.  It must transport the differential,
cycles, boundaries and filling operators together.  Merely approximating
boundary matrices does not suffice: an approximately square-zero matrix is
not a chain complex.

The algebraic tool is the classical homological perturbation lemma
\cite{crainic2004}.  ProveIt's unknot-recognition report~13 already implements
a finite version over $\mathbb F_2$, with a nilpotence bound coming from a
positive grading~\cite{repo_hpl}.  The contribution here is its integration
with quantitative presentations: signed rational formulas, an entirely
finite residual theorem, explicit precision schedules, and a bounded-operator
interpretation.  The underlying perturbation lemma is established
mathematics.  The new interface avoids repeatedly rebuilding a homological
model when its differential changes.

\subsection{Data, signs and an exact algebraic theorem}

A bounded chain complex is a finite family of vector spaces $C_j$ with maps
$d_j:C_j\to C_{j-1}$ satisfying $d_{j-1}d_j=0$.  Write $C$ for their direct
sum and $d$ for the resulting degree $-1$ map.  Initially the spaces can be
finite-dimensional rational coordinate spaces; later they may be Banach
spaces with bounded structural maps.  A smaller complex $(B,b)$ comes with
degree-zero maps $i:B\to C$, $p:C\to B$ and a degree-one map $h:C\to C$.
Our input is the complete list of identities
\begin{align}
 d^2=b^2&=0,& di&=ib,& pd&=bp,& pi&=1_B,\\
 dh+hd&=1_C-ip,& hi&=0,& ph&=0,& h^2&=0.
 \label{eq:hpl-input}
\end{align}
These data are a normalized strong deformation retraction.  Multiplication
of operators means composition, with the rightmost operator acting first.
The sign in $dh+hd=1-ip$ fixes every sign below.  In the convention
$ip=1+dh+hd$ of~\cite{crainic2004}, replace $h$ by $-h$ to obtain ours.

A perturbation is a degree $-1$ map $\delta:C\to C$ for which
\begin{equation}
 (d+\delta)^2=0,
 \quad\text{equivalently}\quad d\delta+\delta d+\delta^2=0.
 \label{eq:hpl-square-zero}
\end{equation}
It is important that this is an exact hypothesis.  Put $d'=d+\delta$.

\begin{theorem}[Signed perturbation with explicit witnesses]
\label{thm:hpl-exact}
Assume \eqref{eq:hpl-input}--\eqref{eq:hpl-square-zero}.  Suppose that
$1+h\delta$ and $1+\delta h$ have two-sided inverses $F,G$, respectively.
Define
\begin{align}
 A&=G\delta=\delta F,&
 b'&=b+pAi,\\
 i'&=i-hAi=Fi,&
 p'&=p-pAh=pG,&
 h'&=h-hAh=Fh=hG.
 \label{eq:hpl-transferred}
\end{align}
Then $(B,b')$ is a complex and
\begin{align}
 (b')^2&=0,& d'i'&=i'b',&p'd'&=b'p',&p'i'&=1_B,\\
 d'h'+h'd'&=1_C-i'p',&h'i'&=0,&p'h'&=0,&(h')^2&=0.
 \label{eq:hpl-output}
\end{align}
If $(h\delta)^L=(\delta h)^L=0$ for a supplied positive integer $L$, the
inverses and all witnesses are finite rational expressions:
\[
 F=\sum_{j=0}^{L-1}(-h\delta)^j,
 \qquad G=\sum_{j=0}^{L-1}(-\delta h)^j.
\]
\end{theorem}

\begin{proof}
The identity $(1+\delta h)\delta=\delta(1+h\delta)$ gives $G\delta=\delta F$.
From the two inverse equations,
\begin{equation}
 A=\delta-\delta hA=\delta-Ah\delta,
 \qquad F=1-hA,\qquad G=1-Ah.
 \label{eq:hpl-resolvent}
\end{equation}
Put $e=ip$, $L_0=1+\delta h$ and $R_0=1+h\delta$.  A key cancellation is
\begin{align*}
 L_0(dA+Ad)R_0
 &=d\delta+\delta d+\delta(hd+dh)\delta\\
 &=-\delta^2+\delta(1-e)\delta=-\delta e\delta.
\end{align*}
Multiplication by the two inverses therefore yields
\begin{equation}
 dA+Ad+AeA=0.
 \label{eq:hpl-maurer}
\end{equation}
This is an algebraic identity, not a convergence argument.  It proves
square-zero of the new differential directly:
\[
 (b+pAi)^2
 =b^2+p(dA+Ad+AeA)i=0.
\]
The two original chain-map identities were used to move $b$ across $p$ and
$i$ in this expansion.

For the inclusion, use $dh=1-e-hd$ and \eqref{eq:hpl-resolvent} to compute
\[
 d'(i-hAi)=ib+eAi+hdAi.
\]
Equation~\eqref{eq:hpl-maurer} changes the last term into
$-hAdi-hAeAi$, which gives $(i-hAi)(b+pAi)$.  Similarly
\[
 (p-pAh)d'=bp+pAe+pAdh
              =(b+pAi)(p-pAh).
\]
Expanding $(p-pAh)(i-hAi)$ leaves $pi=1_B$: every other term contains
$ph$, $hi$ or $h^2$.

For the homotopy identity, first absorb the terms containing a single
$\delta$ using \eqref{eq:hpl-resolvent}.  The full calculation is
\begin{align*}
 d'h'+h'd'
 &=1-e+Ah+hA-dhAh-hAhd\\
 &=1-e+eAh+hAe+h(dA+Ad)h\\
 &=1-e+eAh+hAe-hAeAh\\
 &=1-(i-hAi)(p-pAh).
\end{align*}
The three remaining side conditions follow by expansion; each summand
contains one of $h^2$, $hi$, $ph$.  Finally,
$(1+T)\sum_{j<L}(-T)^j=1-(-T)^L$, on either side.  A supplied
nilpotence exponent therefore supplies the two inverses by finite sums.
\end{proof}

The theorem gives a quotient-free formulation of homology preservation.
If $d'z=0$, then $b'p'z=0$ and
\begin{equation}
 z-i'p'z=d'h'z.
 \label{eq:hpl-cycle-transfer}
\end{equation}
If $b'w=0$, then $d'i'w=0$ and $p'i'w=w$.  If $i'w=d'v$, then
$w=b'p'v$.  These are equations on vectors and explicit witnesses; the
object language need not construct the quotient $\ker d'/\operatorname{im}d'$.

\subsection{Finite residuals without an infinite inverse}

Equip finite rational coordinate spaces with the maximum norm and matrices
with the maximum absolute row-sum norm.  This norm is rational, is computed
by bounded folds, and is submultiplicative.  The same argument applies to
any submultiplicative operator norm.  Let the supplied nonnegative rational
bounds be
\[
 \|d\|\le D,\quad \|h\|\le H,\quad \|\delta\|\le U,
 \quad\|i\|\le I,\quad\|p\|\le P,
 \quad\|h\delta\|,\|\delta h\|\le q<1.
\]
One can use $q=HU$ when that is smaller than one; computing the two products
can give a much sharper certificate.  No inverse, spectrum, or rank is
part of this input.

For $N\ge0$ put
\begin{align}
 A_N&=\delta\sum_{j=0}^{N-1}(-h\delta)^j,
 &R_N&=\delta(-h\delta)^N=(-\delta h)^N\delta,\\
 M_N&=dA_N+A_Nd+A_NipA_N.
 \label{eq:hpl-finite-data}
\end{align}
The empty sum is zero.  Define $b_N,i_N,p_N,h_N$ by the formulas
\eqref{eq:hpl-transferred}, using $A_N$ instead of $A$.

\begin{theorem}[A finite homological residual certificate]
\label{thm:hpl-finite}
Let
\begin{equation}
 K_M=\frac{U\bigl(2(1+q)D+3UIP\bigr)}{(1-q)^2}.
 \label{eq:hpl-K}
\end{equation}
For every $N\ge0$, the following identities hold exactly:
\begin{align}
 p_Ni_N&=1_B,&h_Ni_N&=0,&p_Nh_N&=0,&h_N^2&=0,\label{eq:hpl-finite-normalizations}\\
 b_N^2&=pM_Ni,\label{eq:hpl-residual-first}\\
 d'i_N-i_Nb_N&=R_Ni+hM_Ni,\\
 p_Nd'-b_Np_N&=pR_N+pM_Nh,\\
 d'h_N+h_Nd'-(1-i_Np_N)&=R_Nh+hR_N+hM_Nh.
 \label{eq:hpl-residual-identities}
\end{align}
Moreover $\|M_N\|\le K_Mq^N$, and consequently the four potentially nonzero residuals
in \eqref{eq:hpl-residual-first}--\eqref{eq:hpl-residual-identities} are bounded, in their displayed order,
by
\begin{equation}
 PI K_M q^N,\quad
 I(U+HK_M)q^N,\quad
 P(U+HK_M)q^N,\quad
 (2HU+H^2K_M)q^N.
 \label{eq:hpl-residual-bounds}
\end{equation}
All assertions concern finite arithmetic data.
\end{theorem}

\begin{proof}
The finite geometric identity gives
\[
 \delta-A_N-\delta hA_N
 =\delta-A_N-A_Nh\delta=R_N,
 \qquad \|R_N\|\le Uq^N.
\]
Repeating the expansions in Theorem~\ref{thm:hpl-exact}, retaining these
two remainder terms, gives every identity in
\eqref{eq:hpl-finite-normalizations}--\eqref{eq:hpl-residual-identities}.  In particular the new differential's
square is a computed defect; it is not silently assumed to vanish.

Here is a direct finite bound on that defect's common source.  With
$e=ip$, $L_0=1+\delta h$, $R_0=1+h\delta$, multiplication gives
\begin{align*}
 L_0M_NR_0={}&-L_0dR_N-R_NdR_0
              -\delta eR_N-R_Ne\delta+R_NeR_N.
\end{align*}
For any matrix $X$, the triangle inequality gives
\[
 (1-q)\|X\|\le\|L_0X\|,
 \qquad (1-q)\|X\|\le\|XR_0\|.
\]
Thus $(1-q)^2\|M_N\|\le\|L_0M_NR_0\|$, without constructing an
inverse.  Use $\|e\|\le IP$, $\|L_0\|,\|R_0\|\le1+q$ and
$q^N\le1$ to obtain
\[
 (1-q)^2\|M_N\|
 \le \bigl(2(1+q)D+3UIP\bigr)Uq^N.
\]
Dividing by the explicitly positive rational $(1-q)^2$ proves the claim.
Submultiplicativity applied to the exact residual identities gives the four
remaining bounds.
\end{proof}

The constants in this finite theorem are deliberately conservative.  Its
advantage is that the arithmetic proof never quantifies over an infinite
operator or uses completeness.  A finite truncation generally has
$b_N^2\ne0$; until the residual is exactly zero it is an approximation to a
complex, not a new exact complex.

\subsection{Primitive-recursive clocks and the continuous interpretation}

\begin{lemma}[A rational geometric clock]
\label{lem:hpl-clock}
For $0\le q<1$ and rational $K\ge1$, let
$s=\lceil(1-q)^{-1}\rceil$.  Choose an integer $c$ with $K\le2^c$,
for example the bit length of the numerator in a positive rational code
for $K$.  Then
\[
 N=s(k+c)\quad\Longrightarrow\quad Kq^N\le2^{-k}.
\]
Every operation in this schedule is primitive recursive on rational codes.
\end{lemma}

\begin{proof}
We have $q\le1-1/s$.  The polynomial inequality
$(1-t)^n(1+nt)\le1$ for $0\le t\le1$ follows by induction: multiplying
the next factor uses
$(1-t)(1+(n+1)t)=1+nt-(n+1)t^2\le1+nt$.
At $t=1/s$ and $n=s$ this gives $q^s\le1/2$, including $s=1$.
Taking $k+c$ blocks and multiplying by $K\le2^c$ gives the result.
After denominator clearing these are quantifier-free rational inequalities
and induction on natural-number counters.
\end{proof}

Take $K$ to be at least one and at least the four constants in
\eqref{eq:hpl-residual-bounds}.  This supplies a single explicit clock for
all homological equations.  The original coordinate dimensions, the
iteration bound $N$, and the rational operations determine every loop.
There is no search for a stage at which an unquantified limit is sufficiently
close.

\begin{theorem}[Bounded-operator completion and stability]
\label{thm:hpl-banach}
Suppose the spaces in \eqref{eq:hpl-input} are Banach spaces, all displayed
maps are bounded, and the same rational norm bounds and $q<1$ are supplied.
Then the series defining $F,G$ converge, Theorem~\ref{thm:hpl-exact} applies,
and
\begin{align}
 \|A-A_N\|&\le\frac{Uq^N}{1-q},\\
 \|i'-i_N\|&\le\frac{Iq^{N+1}}{1-q},&
 \|p'-p_N\|&\le\frac{Pq^{N+1}}{1-q},\\
 \|h'-h_N\|&\le\frac{Hq^{N+1}}{1-q},&
 \|b'-b_N\|&\le\frac{PIUq^N}{1-q}.
 \label{eq:hpl-tail}
\end{align}
Also $\|i'\|\le I/(1-q)$, $\|p'\|\le P/(1-q)$ and
$\|h'\|\le H/(1-q)$.

For two perturbations $\delta_1,\delta_2$ of the same original contraction,
suppose both have the product bound $q$ and
$\|\delta_1-\delta_2\|\le\eta$.  Then
\begin{equation}
 \|A_1-A_2\|\le\frac{\eta}{(1-q)^2}.
 \label{eq:hpl-stability}
\end{equation}
The corresponding differences of $b',i',p',h'$ are bounded by this
quantity multiplied by $PI,HI,PH,H^2$, respectively.
\end{theorem}

\begin{proof}
For any finite $M$ the tails of the geometric sums have norm at most
$q^N\sum_{j<M}q^j\le q^N/(1-q)$.  Thus the explicit clock above makes
the sums Cauchy.  Completeness supplies their limits in the bounded-operator
space.  The finite inverse defects are $(-h\delta)^N$ and
$(-\delta h)^N$, tending to zero; continuity of multiplication gives the
two-sided inverse equations.  This proves existence with a rate, rather
than appealing to a qualitative inverse theorem.

The first tail bound follows on multiplying by $\delta$.  The identity
$1-hA_N=\sum_{j=0}^{N}(-h\delta)^j$ shows why the inclusion and homotopy
bounds have exponent $N+1$; the projection has the analogous identity.
Multiplying by $p,i,h$ gives \eqref{eq:hpl-tail}.  Summing the geometric
majorant gives the norm bounds.

Write $G_j=(1+\delta_jh)^{-1}$ and $F_j=(1+h\delta_j)^{-1}$.
The resolvent identity is
\[
 A_1-A_2=G_1(\delta_1-\delta_2)F_2.
\]
Indeed $G_1-G_2=-G_1(\delta_1-\delta_2)hG_2$; substitute this in
$G_1\delta_1-G_2\delta_2$ and use $1-hA_2=F_2$.
The two inverse norm bounds are $1/(1-q)$, proving
\eqref{eq:hpl-stability}.  The remaining assertions follow from
\eqref{eq:hpl-transferred}.
\end{proof}

This Banach statement is the semantic interpretation.  Its iPRA content is
the family of finite tail and residual inequalities, represented by rational
codes and proved with the stated clocks.  To use infinite-dimensional
operators computationally, one additionally supplies a rational operator
presentation with primitive-recursive composition, evaluation and certified
norm bounds, or a fixed algorithm with proofs of these properties.
An arbitrary bounded operator is not automatically such an input.
Unbounded differentials require separate invariant-domain and graph-norm
arguments and are outside this theorem.

\subsection{Non-acyclic overlaps and quantitative fillings}

\begin{corollary}[A global model from supplied local models]
\label{cor:hpl-bicomplex}
Let $C_{a,j}$, $0\le a\le r$, be a finite-column double complex.  Its
vertical differential $d$ lowers $j$, its horizontal differential $\delta$
lowers $a$, and
$d^2=\delta^2=d\delta+\delta d=0$.  Suppose each column carries data
\eqref{eq:hpl-input} onto a supplied smaller complex, with $i,p,h$
preserving the column number.  Then their direct sum retracts the total
complex onto the smaller total space with differential
\begin{equation}
 b'=b+\sum_{j=0}^{r}(-1)^j p\delta(h\delta)^ji.
 \label{eq:hpl-global}
\end{equation}
The retraction maps and all cycle and boundary witnesses are finite
primitive-recursive expressions.  No local acyclicity assumption is needed.
\end{corollary}

\begin{proof}
Both $h\delta$ and $\delta h$ lower the column number.  Their $(r+1)$st
powers therefore vanish: on every basis block, no nonzero target column
remains.  Apply the nilpotent case of Theorem~\ref{thm:hpl-exact} with
$L=r+1$.  The displayed formula is $b+pAi$ with that finite inverse.
All chain identities and the explicit equations
\eqref{eq:hpl-cycle-transfer} follow from the same theorem.  A signed
totalization of a geometric double complex must supply the stated
anticommutation identity; it is not inferred from approximate matrices.
\end{proof}

For a finite cover this allows each intersection to retain its own explicit
homological model.  The terms with $j>0$ account for successive excursions
through eliminated pieces.  When intersections have nontrivial homology,
the result is generally larger than the ordinary nerve.  Thus the earlier
acyclic-overlap restriction is replaced by computable local models, with
the corrected global differential recorded explicitly.  Constructing those
local models is a remaining geometric task for each application.

Here is a precise finite-cover instance, including its comparison with the
original space.  Let a finite ordered simplicial complex $K$ be the union
of subcomplexes $K_1,\ldots,K_m$.  In bidegree $(a,j)$ take the direct sum
of rational simplicial $j$-chains on $K_{i_0}\cap\cdots\cap K_{i_a}$,
over strictly increasing index tuples.  The horizontal map deletes a cover
index with the alternating sign.  The vertical map is $(-1)^a$ times the
simplicial boundary.  The two maps anticommute because deleting a cover
index changes the parity of $a$, while restriction commutes with the
ordinary simplicial boundary.

For each simplex $\sigma$ of $K$, the nonempty set of cover indices
containing it has a least member $u(\sigma)$.  Its horizontal row is the
chain complex of a full simplex on those indices.  Coning an index tuple
to $u(\sigma)$ contracts this row onto one copy of $\sigma$: insert
$u(\sigma)$ into the ordered tuple with the insertion sign, and return
zero if it was already present.  In the boundary of this cone, the term
deleting the inserted vertex is the original tuple; every other deletion
cancels its counterpart obtained by coning after deletion.  In degree zero
the remaining term is the chosen vertex.  This proves
$\delta h_0+h_0\delta=1-i_0p_0$.  Also $h_0^2=h_0i_0=p_0h_0=0$,
since insertion is performed at most once and raises the horizontal degree.
Here $p_0$ forgets the singleton cover index and is zero in higher
horizontal degree, and $i_0$ chooses $u(\sigma)$.

Apply Theorem~\ref{thm:hpl-exact} to this horizontal contraction, with the
vertical boundary as perturbation.  Each product of that boundary with
$h_0$ lowers the simplicial degree, so its $(\dim K+1)$st power is zero
for a nonempty complex.  Every positive correction term in the transferred
small differential raises the horizontal degree before applying $p_0$,
and is therefore killed.  The surviving term $p_0di_0$ is exactly the
simplicial boundary of $K$.  Thus the total double complex contracts onto
$C_*(K;\Q)$ by explicit finite formulas.  Applying
Corollary~\ref{cor:hpl-bicomplex} in the other direction, using the supplied
intersection models, gives an explicit homology comparison between those
updated small models and $K$.  The empty complex is handled by the zero
complex.  No appeal to a spectral sequence or an unconstructed comparison
map is needed.  This also clarifies the limitation: the model may be large
when the intersections have large homology, and no uniform speed advantage
over direct elimination is asserted.

The continuous version also gives useful estimates on actual fillings.
If $d'z=0$ and a vector $s$ in the small model satisfies $b's=p'z$, then
\[
 v=h'z+i's,\qquad d'v=z,
 \qquad
 \|v\|\le\frac{H\|z\|+I\|s\|}{1-q}.
\]
The equation follows by adding the two chain identities.  In particular,
when $B=0$, every cycle has the canonical filling $h'z$ with norm at most
$H\|z\|/(1-q)$.  Conversely, a supplied linear functional $\lambda$ on
$B$ with $\lambda b'=0$ and $\lambda p'z\ne0$ proves that $z$ cannot be a
boundary.  A positive rational lower bound for $|\lambda p'z|$ makes this
obstruction robust under certified approximation.  These are explicit
filling or obstruction certificates, without an undecidable test of equality
of arbitrary real coefficients.

\subsection{A continuous Schur-complement example and a necessary boundary}

Take a two-term complex $C_1=C_0=\Q^2$ and
\[
 d=\begin{pmatrix}1&0\\0&0\end{pmatrix},\qquad
 B_1=B_0=\Q,\quad b=0.
\]
Let $i$ include the second coordinate, $p$ project to it, and let
$h:C_0\to C_1$ be the matrix $\operatorname{diag}(1,0)$, with $h=0$ on
$C_1$.  These satisfy all the normalized contraction identities on the
graded direct sum.  Perturb the only differential by
\[
 \delta=\begin{pmatrix}a&\beta\\\gamma&\tau\end{pmatrix}.
\]
There is no next differential, so square-zero of $d+\delta$ is automatic.
Whenever the supplied margin makes $1+a$ invertible, the transferred
differential is
\begin{equation}
 b'=\tau-\frac{\gamma\beta}{1+a}.
 \label{eq:hpl-schur}
\end{equation}
Indeed $i'$ on degree one is the column
$(-\beta/(1+a),1)^{\mathsf T}$, $p'$ on degree zero is the row
$(-\gamma/(1+a),1)$, and the surviving entry of $h'$ is $1/(1+a)$.
Multiplication by $d+\delta$ gives \eqref{eq:hpl-schur} directly.

For $(a,\beta,\gamma,\tau)=(1/4,1/8,1/5,1/50)$, the small differential
is zero and the product norms are bounded by $q=3/8$.  Changing only
$\tau$ to $1/25$ changes $b'$ to $1/50$.  Thus a small differential
perturbation can destroy homology on the retained model, even though all
perturbation formulas remain well-conditioned.  The theorem preserves
equivalence to the \emph{updated} small complex, not the original Betti
numbers.  Stable Betti numbers require extra information, such as an
appropriate rank margin or vanishing of the transferred perturbation.

The exact-rational executable \texttt{code/homological\_examples.py}
checks this example against an independently computed rational inverse,
verifies all contraction identities, and checks the finite residual bounds.
It also treats $a=0$, where both product squares vanish and the two-term
finite series is exact even without a norm-smallness premise.

\subsection{A precise Lean and object-language development path}

The smallest semantic interface is a structure of rational matrices with
fields $d,b,i,p,h$, shape and degree data, and proofs of
\eqref{eq:hpl-input}.  A second structure adds $\delta$ and the proof of
\eqref{eq:hpl-square-zero}.  Define \texttt{neumann} by primitive recursion
on $N$, and define $A_N,R_N,M_N$ by \eqref{eq:hpl-finite-data}.
The main theorem dependencies are:
\begin{enumerate}
\item \texttt{geometric\_defect}: both finite inverse defects equal the
appropriate $N$th power; prove by induction on $N$.
\item \texttt{residual\_identities}: the seven groups of identities in
\eqref{eq:hpl-finite-normalizations}--\eqref{eq:hpl-residual-identities}; prove by distributivity and the supplied
chain and normalization equations.  This is noncommutative algebra, so a
commutative \texttt{ring} call on whole matrices is inappropriate.
\item \texttt{rowNorm\_mul\_le}: the rational maximum-row-sum norm is
submultiplicative; prove finite triangle bounds, then the bounded maximum
property.  Use it to prove \texttt{finite\_residual\_bound} exactly as in
Theorem~\ref{thm:hpl-finite}.
\item \texttt{block\_clock}: prove Lemma~\ref{lem:hpl-clock} over rational
codes after denominator clearing.  The invariant is the displayed
polynomial inequality, not a real logarithm.
\item \texttt{filtered\_nilpotent}: a product that lowers a bounded column
index more than $r$ times is zero.  Instantiate \texttt{exact\_hpl} and
\texttt{cycle\_transfer} for the double complex.
\end{enumerate}

For the first-order endpoint, encode an $m$ by $n$ matrix as its dimensions
and a finite list of signed numerator/positive-denominator pairs.  Define
dimension checks, entries, products, row sums, maximums and equality tests
as fixed primitive-recursive functions, with total default values on
malformed codes.  Guard every theorem by a Boolean validity test.  The
universally closed assertion is
\[
 \operatorname{ValidContraction}(w)=1\ \wedge\
 \operatorname{NormGuards}(w,q)=1
 \ \Longrightarrow\
 \operatorname{ResidualBounds}(w,N,q)=1.
\]
Bounded array assertions are represented by Boolean folds.  Their prefix
identities and norm inequalities use quantifier-free induction; the
denominator margin is explicit.  A proof of this sentence in
\texttt{Prf PRAAx}, followed by the repository's established bridge, is the
formal iPRA endpoint.  Ordinary Lean proofs of matrix identities, the
Banach-space interpretation, and exact-rational executions support distinct
claims.  This chapter supplies the complete mathematical argument and
construction; it does not claim that this full object derivation has already
been implemented.
